% Research addendum: not yet a proven theorem about actual Syracuse orbits.
% The implication is elementary; hypothesis (SL2) remains OPEN.
% Intended location: docs/research/. Do not input into main.tex before review.
\subsection*{A dyadic local-time criterion for the first-passage tail}
\label{subsec:shell-local-time}
Retain the setting and notation of Proposition~\ref{prop:return-count-criterion}.
For an odd starting value $m$, let $N_{X,j}(m)$ be the number of occurrences,
counted with multiplicity, in the dyadic shell
\[
 I_{X,j}=\bigl(X2^{-(j+1)},X2^{-j}\bigr]\cap W_X
\]
among exactly the Syracuse blocks used to define $\Phi_X(m)$. The shells
are disjoint, and only finitely many are nonempty for each fixed $X$.
The elementary inequality
\begin{equation}\label{eq:shell-pointwise}
 0\le\frac{\Phi_X(m)}X\le
 \sum_{j\ge0}2^{-j}N_{X,j}(m)
\end{equation}
retains the size information lost in the unweighted return count $V_X$.

\begin{proposition}[Dyadic local-time second-moment criterion]
\label{prop:shell-moment-criterion}
Assume the existence of the logarithmic means used in
Proposition~\ref{prop:return-count-criterion}. If a constant $B<\infty$
satisfies, for all sufficiently large $X$ and every $j\ge0$,
\begin{equation}\label{eq:shell-moment-open}
 \mathbb E_{\log}\bigl[N_{X,j}(w^\flat)^2\bigr]\le B,
\end{equation}
then
\begin{equation}\label{eq:shell-tail-bound}
 \mathbb E_{\log}\bigl[(\Phi_X(w^\flat)/X)^2\bigr]\le4B,
 \qquad
 0\le T_X\le 4B\,X(\log X)^{-\varepsilon_0/2}=o(X).
\end{equation}
This is a conditional criterion. Equation~\eqref{eq:shell-moment-open}
for actual Syracuse orbits is not claimed here.
\end{proposition}
\begin{proof}
For every fixed $X$ the sum in~\eqref{eq:shell-pointwise} is finite.
For $j,k\ge0$, Cauchy--Schwarz gives
\[
 \mathbb E_{\log}[N_{X,j}N_{X,k}]
 \le \sqrt{\mathbb E_{\log}[N_{X,j}^2]}
      \sqrt{\mathbb E_{\log}[N_{X,k}^2]}\le B.
\]
Expanding the square in~\eqref{eq:shell-pointwise} and using
$\sum_{j\ge0}2^{-j}=2$ therefore yields
\[
 \mathbb E_{\log}\left[(\Phi_X/X)^2\right]
 \le B\bigl(\sum_{j\ge0}2^{-j}\bigr)^2=4B.
\]
For $y\ge0$ and $H>0$, $(y-H)_+\le y^2/H$. Apply this with
$y=\Phi_X$ and $H=H_X=X(\log X)^{\varepsilon_0/2}$ and take
logarithmic means to obtain~\eqref{eq:shell-tail-bound}.
\end{proof}

\paragraph{Why this is preferable to a uniform moment of $V_X$.}
In the ideal renewal model of Lemma~\ref{lem:model}, the expected
\emph{unweighted} total number of points with logarithms in
$((1-\beta)\log X,\log X]$ is
\[
 \gamma_0\beta\log X+O(1),\qquad
 \gamma_0=2/\log(4/3),
\]
by the renewal estimate $V(t)=\gamma_0t+O(1)$, averaged over the
uniform logarithmic starting interval. Thus the full visit count grows
in this model while each fixed-width logarithmic shell has expected
count $O(1)$. This model statement does not, by itself, prove the
corresponding growth for actual Syracuse orbits. It warns against
using a scale-independent moment bound for the raw total count as the
principal conjecture.

\paragraph{The next open assertion.}
The meaningful next theorem to seek is the actual-orbit bound
\eqref{eq:shell-moment-open}, or a weighted substitute sufficient to
bound the right-hand side of~\eqref{eq:shell-pointwise} in $L^2$.
The ideal-model local visit bounds in Lemma~\ref{lem:model} do not
transfer to the exceptional coupling events without an additional
uniform-integrability estimate. Moreover this criterion does not
control the separate residual term $R_X$.
